\subsection{Representative Comparison of $N_{\textup{sp}}=8000$ and $N_{\textup{sp}}=2000$ Runs}\label{app:nsp8000_vs_2000}

To provide a matched comparison of computational burden, we contrasted aligned runs with $N_{\textup{sp}}=8000$ against the standardized runs with $N_{\textup{sp}}=2000$ using one representative dependence setting per scenario, namely $\rho=0$. Both sets of runs used the same adaptive Monte Carlo thresholds, $\varepsilon_{\tau}=5\times 10^{-4}$ and $\varepsilon_{\xi}=10^{-4}$, and all computations used 20 cores.

Across all four scenarios, reducing $N_{\textup{sp}}$ from 8000 to 2000 substantially decreased peak worker memory usage from approximately 8--9~GB to about 0.8--0.9~GB. Under this matched-threshold comparison, the $N_{\textup{sp}}=2000$ runs were also faster in all four representative settings. Among the settings that converged for both super-population sizes, the largest time savings were seen in Scenarios~S2 and S4, whereas the difference was more modest in Scenario~S1. The runtime comparison for Scenario~S3 should be interpreted more cautiously: although the $N_{\textup{sp}}=2000$ run finished sooner, it reached $B_{\max}=3000$, while the aligned $N_{\textup{sp}}=8000$ run converged at $B_{\text{final}}=850$. Thus, the smaller super-population size is clearly attractive from a memory perspective and is often faster in practice, but the survival-plus-continuous setting remains the most sensitive to the reduction in $N_{\textup{sp}}$.

\begin{table}[ht!]
\centering
\caption{Representative computational comparison between aligned $N_{\textup{sp}}=8000$ runs and standardized $N_{\textup{sp}}=2000$ runs at $\rho=0$.}
\label{tab:nsp8000_vs_2000_rho0}
\setlength{\tabcolsep}{3.5pt}
\footnotesize
\begin{tabular}{@{}l c ccc ccc@{}}
\toprule
& & \multicolumn{3}{c}{$N_{\textup{sp}}=8000$} & \multicolumn{3}{c}{$N_{\textup{sp}}=2000$} \\
\cmidrule(lr){3-5} \cmidrule(lr){6-8}
\textbf{Scenario} & $\rho$ & $B_{\text{final}}$ & Time (min) & Worker max (GB) & $B_{\text{final}}$ & Time (min) & Worker max (GB) \\
\midrule
S1: Cont + Cont & 0 & 250 & 12.3 & 7.80 & 900 & 6.5 & 0.80 \\
S2: Cont + Bin & 0 & 350 & 17.4 & 7.80 & 1350 & 7.5 & 0.80 \\
S3: Surv + Cont & 0 & 850 & 45.6 & 8.75 & 3000 & 14.3 & 0.91 \\
S4: Bin + Cont & 0 & 350 & 17.8 & 7.80 & 1400 & 6.4 & 0.81 \\
\bottomrule
\end{tabular}
\parbox{0.92\textwidth}{\footnotesize \textit{Note:} This table uses one representative setting per scenario, namely $\rho=0$, and all runs used 20 cores. The aligned $N_{\textup{sp}}=8000$ diagnostics are taken from \texttt{Diagnostics.Nsp8000.Aligned.Rho0.csv}, which was generated with the same adaptive Monte Carlo thresholds used in the standardized runs, namely $\varepsilon_{\tau}=5\times 10^{-4}$ and $\varepsilon_{\xi}=10^{-4}$. The $N_{\textup{sp}}=2000$ runs are taken from \texttt{Summary.Scenario1.Nsp2000.csv}, \texttt{Summary.Scenario2.Nsp2000.csv}, \texttt{Summary.Scenario3.Nsp2000.csv}, and \texttt{Summary.Scenario4.Nsp2000.csv}. Under this matched-threshold comparison, the smaller super-population size still greatly reduces worker memory and is faster in all four representative settings. However, the apparent runtime advantage for the survival-plus-continuous setting (S3) should be interpreted cautiously because the $N_{\textup{sp}}=2000$ run reached $B_{\max}$ rather than full convergence.}
\end{table}

